%% introduction.tex
\section{Introduction}\label{sec:intro}

Biotechnology depends on predicting and directing cellular metabolism.
Metabolic engineering redirects flux toward a product. Bioprocess
development selects operating conditions that sustain production. Strain
design identifies genetic changes that alter the phenotype. Each task
requires an estimate of metabolic reaction rates under a specified condition
\cite{Bordbar2014}. These rates are difficult to measure directly. Genome
sequences and reaction annotations provide a network structure, but
reconstructions remain uncertain in reaction membership, directionality,
compartments, cofactors, and biomass composition. Fluxes, kinetic parameters,
and metabolite concentrations are even less complete at genome scale. Models
must therefore estimate phenotype from incomplete measurements. Bioprocess
models use several levels of biological detail. Unstructured
models track biomass, extracellular substrates, and products with lumped
growth, uptake, and production rates. Monod growth, inhibition, maintenance,
and growth-associated production laws support batch, fed-batch, and
continuous reactor analysis, but do not resolve intracellular fluxes
\cite{Esener1983}. Structured kinetic models add intracellular pools,
enzymes, or cellular machinery. Early single-cell models from Shuler and
coworkers coupled intracellular metabolism to extracellular nutrients in
microbial and mammalian cells \cite{Domach1984,Wu1992}. Reaction-level
kinetic models assign mass-action, Michaelis--Menten, or power-law rates to
individual reactions and integrate the resulting differential equations
\cite{MichaelisMenten1913,Savageau1976,SavageauVoit1987,Vadhin2026}. These
models describe transient intracellular behavior, but require more kinetic
parameters as their resolution increases.

Constraint-based models use a different closure. They impose mass
conservation on the reaction network and require fewer explicit kinetic
parameters. This approach can be applied to genome-scale reconstructions
\cite{Bordbar2014,AllenPalsson2003}. Flux balance analysis (FBA) formulates
the calculation as a linear program \cite{Orth2010,Heirendt2019}.
Conservation alone usually permits many flux distributions. Reaction bounds
restrict this feasible set, and an objective selects an optimum. Bounds are
a major route by which condition-specific thermodynamic, kinetic,
expression, and regulatory data enter the model. This chapter examines how
each type of data changes a bound and its predicted fluxes. The chapter first
derives the flux constraint from a mole balance, including
the growth-dilution balance $\mathbf S\hat{\mathbf v}=\mu\mathbf x$, where
$\mathbf S$ is the stoichiometric matrix, $\hat{\mathbf v}$ is the
intracellular flux vector, $\mu$ is the growth rate, and $\mathbf x$ is the
intracellular concentration vector. It then
factors each reaction bound into terms for reversibility, enzyme capacity,
substrate saturation, enzyme abundance, and regulation. Two small examples
show how these terms affect a flux prediction. The first uses public data to
set bounds in a urea-cycle network. The second transfers expression and
metabolite states from a kinetic feedback model into an FBA bound. Two larger
studies illustrate related methods in dynamic cell-free protein synthesis
\cite{Dai2018,Vilkhovoy2023} and model-guided strain design \cite{Wayman2019}.
Strain design also requires a discrete search over candidate interventions,
often formulated as a bilevel optimization problem \cite{Burgard2003}.
